\documentclass[11pt]{article}
\usepackage[margin=1.1in]{geometry}
\usepackage{amsmath,amssymb,amsthm}
\usepackage{booktabs}
\usepackage{algorithm}
\usepackage{algpseudocode}
\usepackage[hidelinks]{hyperref}

\newtheorem{lemma}{Lemma}
\newtheorem{proposition}{Proposition}
\newtheorem{remark}{Remark}
\newcommand{\li}{\operatorname{li}}
\newcommand{\erfc}{\operatorname{erfc}}
\newcommand{\eps}{\varepsilon}

\title{An analytic algorithm for the sum of primes below $x$}
\author{Raghav Sharma\thanks{Manuscript prepared with substantial AI assistance
(Anthropic Claude); all derivations and computations were machine-verified as
described in Section~\ref{sec:results}. \textbf{Draft — not yet submitted;
rigorous error bounds (Section~\ref{sec:rigor}) remain to be completed.}}}
\date{August 2026}

\begin{document}
\maketitle

\begin{abstract}
We describe an algorithm that computes $S(x)=\sum_{p\le x}p$, the sum of all
primes up to $x$, from the nontrivial zeros of the Riemann zeta function.
The method combines three ingredients: a logarithmic-Gaussian smoothing whose
Mellin transform is exactly Gaussian, so that each zero's contribution is
damped by $e^{-\gamma^2\eps^2/2}$; an integral representation
$1/\log n=\int_0^A n^{-\sigma}\,d\sigma + n^{-A}/\log n$ that removes the
logarithmic weight from the smoothed Chebyshev-type sum, with a prime-zeta
remainder; and exact sieve corrections confined to a narrow window around $x$
and to prime powers. A floating-point implementation recovers the exact
integer value of $S(x)$ for $x$ up to $10^9$ using only the first $2{,}000$
zeros of $\zeta$, e.g.\ $S(10^9)=24{,}739{,}512{,}092{,}254{,}535$ with
residual $6\times10^{-3}$. The structure parallels the Lagarias--Odlyzko
analytic method for $\pi(x)$ as implemented rigorously by Platt; to the best
of our knowledge no analytic computation of \emph{prime sums} has been
published. We state heuristic error bounds, report a cautionary numerical
phenomenon (coherent double-precision rounding in windowed corrections that
masquerades as genuine zero-frequency oscillation), and outline the path to a
rigorous interval-arithmetic version.
\end{abstract}

\section{Introduction}\label{sec:intro}

Computing $\pi(x)$, the number of primes up to $x$, without enumerating
primes is a classical subject: the combinatorial line runs from Legendre and
Meissel through Lehmer to Lagarias--Miller--Odlyzko and
Del\'eglise--Rivat~\cite{DR96}, with record computations by Gourdon, Oliveira
e Silva, Staple~\cite{Staple15}, and Walisch's \texttt{primecount}. The
analytic line, proposed by Lagarias and Odlyzko~\cite{LO87}, evaluates
$\pi(x)$ from the zeros of $\zeta(s)$ in $O(x^{1/2+\eps})$ time; Platt's
rigorous implementation~\cite{Platt15} computed $\pi(10^{24})$
unconditionally.

For the \emph{sum} of primes $S(x)=\sum_{p\le x}p$ the combinatorial theory
transfers: the $O(x^{3/4})$ dynamic-programming method was posted by the
Project Euler user Lucy\_Hedgehog in 2013~\cite{Lucy13}, Fenwick-tree
variants reach $\tilde O(x^{2/3})$~\cite{Broxey23}, and Walisch's
\texttt{primesum}~\cite{Walisch} adapts Del\'eglise--Rivat to sums with
record values past $10^{20}$. The analytic line, however, appears never to
have been carried out for prime sums: we are not aware of any published
explicit-formula computation of $S(x)$. This paper fills that gap at the
level of a working algorithm and verified implementation.

\subsection*{Contributions}
\begin{enumerate}
  \item An explicit formula for Gaussian-smoothed, polynomially weighted
    Chebyshev sums $\sum_n\Lambda(n)\,n^{1-\sigma}c(n)$
    (Proposition~\ref{prop:explicit}), in a form whose zero terms are damped
    by $e^{-\gamma^2\eps^2/2}$ with all transforms in closed form.
  \item A mechanism for removing the $1/\log$ weight
    (Proposition~\ref{prop:logstrip}): integrating the formula over the
    weight exponent $\sigma\in[0,A]$, with the $\zeta'/\zeta$ pole at
    $\sigma=2$ cancelling the main term's pole exactly, and a rapidly
    convergent prime-zeta remainder closing the truncation. This is the
    step with no direct analogue in the $\pi(x)$ literature, where
    Riemann-style $\li$ terms are used instead.
  \item An assembled algorithm (Section~\ref{sec:algorithm}) computing
    $S(x)$ exactly: analytic main computation plus exact sieve corrections
    over a window of relative width $O(\eps)$, with an explicit error
    budget; heuristic cost $O(x^{1/2+\eps})$ when the number of zeros is
    chosen as $N\sim x^{1/2}$ (Section~\ref{sec:complexity}).
  \item A verified implementation: exact integer recovery of $S(x)$ for
    $x=10^6,\dots,10^9$ from $2{,}000$ zeros
    (Table~\ref{tab:results}), cross-checked against independent
    combinatorial computations.
  \item A numerical-analysis caveat of independent interest
    (Section~\ref{sec:rounding}): windowed corrections evaluated in IEEE
    double precision develop a \emph{coherent} rounding bias of size
    $\approx x^2\cdot2^{-52}$ that oscillates at the frequencies of the
    zeta zeros --- indistinguishable in signature from a missing term of the
    explicit formula. We document the effect and its remedy.
\end{enumerate}

\section{Setup}\label{sec:setup}

Throughout, $x\ge2$ is the (integer) target, $\eps>0$ a smoothing parameter,
$L=\log x$, and $\rho=\beta+i\gamma$ runs over nontrivial zeros of $\zeta$;
for the zeros we use ($\gamma\le2515.3$) the Riemann hypothesis is verified,
so $\rho=\tfrac12+i\gamma$, but the formulas hold with $\beta$ in place of
$\tfrac12$. Define the logarithmic-Gaussian cutoff
\begin{equation}\label{eq:cutoff}
  c(t)\;=\;\tfrac12\,\erfc\!\Bigl(\frac{\log(t/x)}{\sqrt2\,\eps}\Bigr),
  \qquad t>0,
\end{equation}
which decreases from $1$ to $0$ across a window of relative width
$O(\eps)$ centred at $x$.

\begin{lemma}[Mellin transform]\label{lem:mellin}
For $\operatorname{Re}s>0$,
\[
  M(s)\;:=\;\int_0^\infty c(t)\,t^{s-1}\,dt\;=\;\frac{x^s}{s}\,
  e^{s^2\eps^2/2},
\]
and the right-hand side continues $M$ meromorphically to $\mathbb C$ with a
single simple pole at $s=0$ of residue $1$.
\end{lemma}

\begin{proof}
Substitute $t=xe^u$, integrate by parts, and use the Gaussian moment
generating function:
$\int_{\mathbb R}\tfrac12\erfc(u/\sqrt2\eps)\,e^{su}\,du
 =\tfrac1s\int_{\mathbb R}\frac{e^{-u^2/2\eps^2}}{\sqrt{2\pi}\,\eps}e^{su}du
 =\tfrac1s e^{s^2\eps^2/2}$;
the boundary terms vanish for $\operatorname{Re}s>0$.
\end{proof}

\begin{lemma}[Smoothed Perron]\label{lem:perron}
For any $c_0>0$ and integer $n\ge1$,
\[
  \frac1{2\pi i}\int_{(c_0)} M(s)\,n^{-s}\,ds\;=\;c(n).
\]
\end{lemma}

\begin{proof}
With $\alpha=\log(x/n)$ the integral is
$F(\alpha)=\frac1{2\pi i}\int_{(c_0)}e^{s\alpha+s^2\eps^2/2}\,\frac{ds}{s}$.
Differentiating under the integral and completing the square gives
$F'(\alpha)=\frac{1}{\sqrt{2\pi}\,\eps}e^{-\alpha^2/2\eps^2}$, independent of
$c_0$; and $F(-\infty)=0$ by closing the contour to the right. Hence $F$ is
the Gaussian distribution function, $F(\alpha)=\Phi(\alpha/\eps)=c(n)$.
\end{proof}

\section{A smoothed explicit formula for weighted sums}\label{sec:explicit}

For $\sigma\ge0$ define
\[
  \widetilde\psi_1(\sigma)\;=\;\sum_{n\ge2}\Lambda(n)\,n^{1-\sigma}c(n),
\]
an absolutely convergent sum since $c$ has compact effective support.

\begin{proposition}\label{prop:explicit}
For $\sigma\ge0$, $\sigma\neq2$,
\begin{equation}\label{eq:explicit}
  \widetilde\psi_1(\sigma)\;=\;M(2-\sigma)\;-\;\sum_{\rho}M(\rho+1-\sigma)
  \;-\;\frac{\zeta'}{\zeta}(\sigma-1)\;-\;\sum_{k\ge1}M(1-2k-\sigma),
\end{equation}
the sum over $\rho$ (with multiplicity, over all $\gamma\in\mathbb R$)
converging absolutely at Gaussian speed
$|M(\rho+1-\sigma)|\ll x^{\beta+1-\sigma}e^{-\gamma^2\eps^2/2}/|\gamma|$.
\end{proposition}

\begin{proof}[Proof sketch]
By Lemma~\ref{lem:perron} and absolute convergence,
$\widetilde\psi_1(\sigma)
 =\frac1{2\pi i}\int_{(c_0)}\bigl[-\tfrac{\zeta'}{\zeta}(s+\sigma-1)\bigr]
 M(s)\,ds$ with $c_0>2-\sigma$. Shift the contour to
$\operatorname{Re}s=-\Omega$. The integrand's poles are: $s=2-\sigma$ (pole
of $\zeta$; residue $M(2-\sigma)$), $s=\rho+1-\sigma$ (residue
$-M(\rho+1-\sigma)$), $s=1-2k-\sigma$ (trivial zeros), and $s=0$ (pole of
$M$; residue $-\tfrac{\zeta'}{\zeta}(\sigma-1)$, using
Lemma~\ref{lem:mellin}). Horizontal segments at height $\pm T$ vanish as
$T\to\infty$ because $M$ decays like $e^{-T^2\eps^2/2}$ against the
polylogarithmic growth of $\zeta'/\zeta$; the left vertical integral tends
to $0$ for $\Omega\to\infty$ chosen along, e.g., $\Omega_j=2j+\tfrac12$,
since $|x^{-\Omega}e^{\Omega^2\eps^2/2}|$ is minimized far beyond all scales
used ($\Omega\ll L/\eps^2$). A rigorous version requires the standard
quantitative forms of these estimates; see Section~\ref{sec:rigor}.
\end{proof}

\section{Removing the logarithmic weight}\label{sec:logstrip}

The object we can assemble $S(x)$ from is
\[
  \widetilde T\;=\;\sum_{p^k}\frac{p^k}{k}\,c(p^k)
  \;=\;\sum_{n\ge2}\frac{\Lambda(n)}{\log n}\,n\,c(n).
\]

\begin{proposition}\label{prop:logstrip}
Fix $A>1$ (we use $A=4$). Then
\begin{equation}\label{eq:logstrip}
  \widetilde T\;=\;\int_0^A\widetilde\psi_1(\sigma)\,d\sigma\;+\;R_A,
  \qquad
  R_A=\sum_{p^k}\frac1k\,p^{-(A-1)k}\,c(p^k),
\end{equation}
and, inserting \eqref{eq:explicit} and integrating term by term
(justified for the zero and trivial sums by uniform Gaussian domination on
$[0,A]$),
\begin{equation}\label{eq:Tfinal}
  \widetilde T=\int_0^A\Bigl[M(2-\sigma)-\frac{\zeta'}{\zeta}(\sigma-1)\Bigr]
  d\sigma
  \;-\;\sum_\rho\int_0^A M(\rho+1-\sigma)\,d\sigma
  \;-\;\sum_{k\ge1}\int_0^A M(1-2k-\sigma)\,d\sigma
  \;+\;R_A .
\end{equation}
The first integrand extends continuously across $\sigma=2$ with value
$\log x-\gamma_E+O(\eps^2)$ there: the simple pole of
$-\zeta'/\zeta(\sigma-1)$ cancels the pole of $M(2-\sigma)$ exactly.
\end{proposition}

\begin{proof}
Equation~\eqref{eq:logstrip} is the identity
$\frac1{\log n}=\int_0^A n^{-\sigma}\,d\sigma+\frac{n^{-A}}{\log n}$
applied under the (absolutely convergent) sum; note
$\Lambda(p^k)/\log p^k=1/k$. For the cancellation at $\sigma=2$, write
$w=2-\sigma$: $M(w)=\frac1w+\log x+O(w)$ (using
Lemma~\ref{lem:mellin}) while
$\frac{\zeta'}{\zeta}(1-w)=\frac1w+\gamma_E+O(w)$.
\end{proof}

\begin{remark}
$R_A$ converges like a prime zeta series: for $A=4$ its $k=1$ part is
$\sum_p p^{-3}c(p)$, and truncating at $p\le10^6$ leaves a tail below
$4\times10^{-14}$. This is what keeps the truncation at finite $A$ cheap:
no enumeration anywhere near $x$ is required.
\end{remark}

\section{The algorithm}\label{sec:algorithm}

From $\widetilde T$ to $S(x)$:
\begin{equation}\label{eq:assembly}
  S(x)\;=\;\widetilde T
  \;-\;\underbrace{\sum_{k\ge2}\frac1k\,p^k c(p^k)}_{\text{prime powers, }
    p\le(xe^{K\eps})^{1/2}}
  \;-\;\underbrace{\sum_{p}\,p\,\bigl(c(p)-[p\le x]\bigr)}_{\text{window
    } |\log(p/x)|\le K\eps}
  \;\;(+\,\text{error}),
\end{equation}
then round to the nearest integer. We take $K=12$
($\tfrac12\erfc(12/\sqrt2)<10^{-32}$, so truncating $c$ outside the window
is negligible against all other budget lines).

\begin{algorithm}
\caption{Analytic $S(x)$}
\begin{algorithmic}[1]
\State choose $\eps$ so that the zero-tail bound
  $2\sum_{\gamma>\gamma_N}x^{3/2}e^{(9/4-\gamma^2)\eps^2/2}/\gamma$
  (zeros beyond the table estimated by the Riemann--von Mangoldt density)
  is below the budget; with $N=2000$ zeros, $\eps\approx8/\gamma_N$
\State $I_{\text{main}}\gets\int_0^A[M(2-\sigma)-\tfrac{\zeta'}{\zeta}
  (\sigma-1)]\,d\sigma$ \Comment{adaptive quadrature, high precision}
\State $I_{\text{zeros}}\gets\sum_{j\le N}2\operatorname{Re}
  \int_0^A M(\rho_j+1-\sigma)\,d\sigma$ \Comment{Gauss--Legendre in
  $u=\sigma L$; integrand smooth, $e^{-u}$-scaled}
\State $I_{\text{triv}}$, $R_A$ \Comment{closed quadratures; sieve to $10^6$}
\State $\widetilde T\gets I_{\text{main}}-I_{\text{zeros}}-I_{\text{triv}}
  +R_A$
\State subtract prime-power and window corrections
  \Comment{sieves to $\sqrt{xe^{K\eps}}$ and over
  $[xe^{-K\eps},xe^{K\eps}]$; \emph{high-precision} $c$, see
  Sec.~\ref{sec:rounding}}
\State \Return nearest integer
\end{algorithmic}
\end{algorithm}

\subsection{Error budget}
With $N$ zeros, table precision $\delta_\gamma$, and working precision
$\delta_f$, the deviation of the computed value from $S(x)$ is bounded
heuristically by the sum of: (i) the zero tail
$O(x^{3/2}e^{-(\gamma_N\eps)^2/2}/(\gamma_N^2\eps^2))$;
(ii) zero-precision noise
$O(x^{3/2}L\,\delta_\gamma\sqrt N)$;
(iii) quadrature and floating-point noise $O(x^{3/2}\delta_f\cdot\text{polylog})$;
(iv) truncation constants below $10^{-13}$ ($R_A$ tail, $K\eps$ cutoffs,
trivial zeros $k\ge3$). All are measured to behave as predicted in
Section~\ref{sec:results}; making (i)--(iii) rigorous is the program of
Section~\ref{sec:rigor}.

\subsection{Complexity (heuristic)}\label{sec:complexity}
The analytic part costs $O(N)$ quadrature-node evaluations (times
$\text{polylog}$ precision factors); the window sieve costs
$O(x\eps\log x)=O((x/\gamma_N)\,\text{polylog})$ since $\eps\asymp
(\log x)^{1/2}/\gamma_N$, and $\gamma_N\sim2\pi N/\log N$. Balancing
$N\asymp x/\gamma_N$ gives $N\asymp x^{1/2}\,\text{polylog}$ and total cost
$O(x^{1/2+\eps})$, the same exponent as Lagarias--Odlyzko for $\pi(x)$
\cite{LO87,Platt15}. Our implementation fixes $N=2000$ and is therefore
window-sieve dominated beyond $x\approx10^8$; scaling $N$ requires only a
larger high-precision zero table (e.g.\ from LMFDB or computed as in
\cite{Platt15}).

\section{Implementation and results}\label{sec:results}

The implementation (\texttt{prototype/analytic\_sum.py} in the accompanying
repository) uses: the first $2{,}000$ zeros computed to 25 significant
digits with mpmath's \texttt{zetazero}; 80-bit extended precision for the
vectorized zero quadrature; mpmath (50 digits) for $I_{\text{main}}$,
assembly, and all window-adjacent cutoff evaluations. Reference values come
from two independent combinatorial implementations (Lucy--Hedgehog
$O(x^{3/4})$ and a Fenwick-hybrid $\tilde O(x^{2/3})$), themselves
cross-checked against OEIS~A046731.

\begin{table}[h]\centering
\begin{tabular}{lllll}
\toprule
$x$ & $S(x)$ & residual & window primes & zeros \\
\midrule
$10^6$ & $37{,}550{,}402{,}023$ & $+2.8\times10^{-6}$ & $4{,}838$ & 2000 \\
$10^7$ & $3{,}203{,}324{,}994{,}356$ & $-2.0\times10^{-5}$ & $44{,}290$ & 2000 \\
$10^8$ & $279{,}209{,}790{,}387{,}276$ & $+6.0\times10^{-4}$ & $414{,}505$ & 2000 \\
$10^9$ & $24{,}739{,}512{,}092{,}254{,}535$ & $+6.0\times10^{-3}$ &
  $3{,}799{,}674$ & 2000 \\
\bottomrule
\end{tabular}
\caption{Exact integer recovery of $S(x)$ from $2{,}000$ zeta zeros.
``Residual'' is the computed value minus the true integer before rounding.
Single-threaded runtimes range from seconds ($10^6$) to $8$ minutes
($10^9$), dominated by high-precision window evaluation.}
\label{tab:results}
\end{table}

The residual grows by roughly one decade per decade of $x$, consistent with
budget line (iii); at fixed working precision the method as configured
would fail to round correctly near $x\approx10^{11}$, and either more
precision, more zeros (shrinking the window), or both restore the margin.

\subsection{A cautionary phenomenon: coherent rounding masquerading as
mathematics}\label{sec:rounding}

An earlier version evaluated the window correction
$\sum_p p\,(c(p)-[p\le x])$ in IEEE double precision. The resulting values
disagreed with the explicit formula by a discrepancy that (a) was
independent of $\eps$, (b) grew like $x^2$, and (c) \emph{oscillated in
$\log x$ at the frequencies of the first zeta zeros} --- precisely the
signature of a missing zero term, which cannot exist for this contour. The
cause: $\log p-\log x$ cancels catastrophically for $p$ in the window, and
the resulting deterministic $\delta c\sim10^{-13}$ errors are weighted by
$p\approx x$ and summed over the $\Theta(x\eps/\log x)$ window primes.
Because the window population itself fluctuates at zero frequencies, the
bias inherits them: an entirely numerical artifact wearing the explicit
formula's clothes. (Verified by recomputing one case wholly in 40-digit
arithmetic: the ``missing term'' collapsed from $-1.2\times10^{-6}$ at
$x=6\times10^4$ to $-3.4\times10^{-10}$, the genuine tail level.) The
remedy costs little: compute $u=\log(p/x)$ as
$\operatorname{log1p}((p-x)/x)$ from exact integers and evaluate $\erfc$
beyond double precision throughout the window. We flag this because any
implementation of windowed analytic prime computations faces it, and the
symptom is maximally misleading.

\section{Toward a rigorous version and records}\label{sec:rigor}

Nothing in the pipeline resists rigor in principle; the required
components exist in the $\pi(x)$ literature. Concretely:
(1) explicit truncation lemmas for Proposition~\ref{prop:explicit}
(quantitative horizontal/vertical contour bounds, as in \cite{Platt15});
(2) an interval-arithmetic implementation with directed rounding through
the quadratures and sieves; (3) a zero database with certified enclosures
(LMFDB provides $10^{11}$ zeros; Platt's methods generate them to any
precision); (4) targets beyond the combinatorial records
($S(x)$ for $x>10^{20}$, cf.~\texttt{primesum}~\cite{Walisch}) require
$N\sim x^{1/2}$ zeros and a parallel window sieve --- the same regime in
which analytic $\pi(x)$ became competitive~\cite{Platt15}. A rigorous
analytic computation of, e.g., $S(10^{22})$, cross-checked against an
independent combinatorial run, would constitute both a new record and an
independent verification methodology, which is where we see the natural
continuation of this work.

\section*{Acknowledgments}
This manuscript, the derivations, and the implementation were produced
with substantial assistance from an AI system (Anthropic's Claude);
the author reviewed and takes responsibility for the content. The zero
table derives from mpmath and cross-checks against Odlyzko's published
tables. All code and data needed to reproduce Table~\ref{tab:results} are
in the accompanying repository.

\begin{thebibliography}{9}
\bibitem{LO87} J.~C. Lagarias and A.~M. Odlyzko, \emph{Computing $\pi(x)$:
an analytic method}, J. Algorithms \textbf{8} (1987), 173--191.
\bibitem{Platt15} D.~J. Platt, \emph{Computing $\pi(x)$ analytically},
Math. Comp. \textbf{84} (2015), 1521--1535. arXiv:1203.5712.
\bibitem{DR96} M. Del\'eglise and J. Rivat, \emph{Computing $\pi(x)$: the
Meissel, Lehmer, Lagarias, Miller, Odlyzko method}, Math. Comp.
\textbf{65} (1996), 235--245.
\bibitem{Staple15} D.~B. Staple, \emph{The combinatorial algorithm for
computing $\pi(x)$}, arXiv:1503.01839 (2015).
\bibitem{Lucy13} Lucy\_Hedgehog, post in Project Euler Problem 10 forum
thread (2013).
\bibitem{Broxey23} G. Broxey, \emph{Lucy's algorithm + Fenwick trees},
blog post (2023),
\url{https://gbroxey.github.io/blog/2023/04/09/lucy-fenwick.html}.
\bibitem{Walisch} K. Walisch, \emph{primesum: sum of the primes below x},
\url{https://github.com/kimwalisch/primesum}.
\bibitem{Odlyzko} A.~M. Odlyzko, \emph{Tables of zeros of the Riemann zeta
function}, \url{https://www-users.cse.umn.edu/~odlyzko/zeta_tables/}.
\bibitem{OEIS} OEIS Foundation, sequence A046731 (sum of primes
$<10^n$), \url{https://oeis.org/A046731}.
\end{thebibliography}

\end{document}
